\section{A paired level and diagram budget}
\label{rec:sec-hybrid}

This proof uses two inductive bounds for the same trace.  The level
budget pays for tuples that lose coordinates on restriction.  The
diagram budget pays for long row and column chains and for distinct
markers outside those chains.  Their common interpolation lemma keeps
the finite starting exponent out of all counting errors.

For $p\ge1$ write
$T_n^\lambda(p)=D_\lambda\Tr|B_n|^{2p}$.
When $T_n^\lambda(p)=0$, we use the extended-real convention
$\log T_n^\lambda(p)=-\infty$.
Fix the constants
\begin{equation}\label{rec:hybrid-constants}
 \begin{gathered}
 c_*=10^{-2},\qquad v_*=10^{-5},\qquad\epsilon=10^{-7},\\
 u_*=1/10,\qquad\Lambda=10/\epsilon^2=10^{15},\qquad
 p_0\ge20\Lambda+10.
 \end{gathered}
\end{equation}
Choose a sufficiently large absolute $d_0$.  Put $\iota_d=0$ for
$d\le d_0$ and
$\iota_d=2d_0^{-1/3}-d^{-1/3}$ otherwise, and set
$(u_d,c_d,v_d)=(u_*,c_*,v_*)+\iota_d(1,1,1)$.
Choose $d_0$ so large that $\iota_d\le v_*/10$.
In particular $\gamma_d=v_d/c_d$ is positive and less than $0.01$.
For a box $x$ of $\lambda$, let $R(x),C(x)$ be the full row and
column lengths, and define
\begin{equation}\label{rec:hybrid-diagram-weight}
 \mathcal D_d(\lambda)=
 \sum_{x\in\lambda}\min\left\{
 c_d\log\frac n{\max(R(x),C(x))},\ v_d\log n\right\}.
\end{equation}

\begin{theorem}\label{rec:hybrid-main}
There is a fixed positive integer $p$ such that, simultaneously for every
integer $d\ge1$ and every partition $\lambda=(n-k,\beta)\vdash n=2^d$
with $\lambda\ne(n)$ and $\lambda'_1\le n/2$, where
$k=n-\lambda_1\ge1$,
\begin{align}
 \log T_n^\lambda(p)&\le
 k\log n\left(u_d-\Lambda
          \left(\frac{\log(n/k)}{\log n}\right)^2\right),
       \label{rec:hybrid-level-budget}\\
 \log T_n^\lambda(p)&\le\mathcal D_d(\lambda).
       \label{rec:hybrid-diagram-budget}
\end{align}
If $\log(n/k)/\log n\ge\epsilon$, the first bound is at most
$-k\log n$.  For all sufficiently large $d$, in the complementary range,
$\mathcal D_d(\lambda)\le\tfrac14\log D_\lambda$.
\end{theorem}

The height condition includes every nonzero block by
Lemma~\ref{lem:annihilation}.  No converse is needed: any zero block
in this height range satisfies both moment inequalities by the
extended-real convention.  The final comparison of the diagram budget
with dimension also covers these zero blocks, because its proof below
uses only the stated height bound and properties of the diagram.

\subsection{Interpolation with a fixed counting exponent}

The two-strip argument below uses classical Schatten three-lines
interpolation; see~\cite[Section 3.1, Theorem 3.1]{SutterBertaTomamichel2017}.
This reference concerns the interpolation principle, while the profile
count and angle inequality below are proved here.

Write the sweep as $B_n=CR$, with $R$ the product of row sweeps
and $C$ the product of column sweeps. Set
\[
 A=(RR^*)^{1/2},\qquad B=(C^*C)^{1/2}.
\]
These positive group-algebra operators still factor over their respective
lines. Within each irreducible block of the corresponding row or column
subgroup algebra choose unitary extensions of the polar factors, so that
$R=AU_R$ and $C=U_CB$. These extensions remain in their subgroup
algebras. Then $CR=U_CBAU_R$,
and $BA$ has the same singular values as the sweep. This identifies the
positive operators used throughout the interpolation argument.

Let a representation $\mathcal H$ contain at least $L_\lambda\ge1$ copies of
$V_\lambda$ in a global $S_n$-invariant subspace $S$, with the same letter denoting its orthogonal projection. In particular $S$ commutes with the row and column operators. Suppose the row and column
profiles give finite families $\{E_\alpha\}$ and $\{F_\beta\}$ of
projections on $\mathcal H$. Within each family the projections are
mutually orthogonal, and
$\sum_\alpha E_\alpha=\sum_\beta F_\beta=I_{\mathcal H}$.
Each projection commutes with $S$; $E_\alpha$ commutes with the row
operator $A$, and $F_\beta$ with the column operator $B$. Set
\[
 J=\#\{(\alpha,\beta):F_\beta E_\alpha S\ne0\},\qquad
 w_{\alpha\beta}=\|F_\beta E_\alpha S\|.
\]
The two projection resolutions and $S\ne0$ imply $J\ge1$.
Then, for $p\ge p_0>1$,
\begin{equation}\label{rec:hybrid-two-strip}
 T_n^\lambda(p)\le
 \frac{D_\lambda}{L_\lambda}J^{2p_0}
 \max_{\alpha,\beta}
 w_{\alpha\beta}^{2(p_0-1)}
 \Tr_{\mathcal H}(B_\beta^{2p}A_\alpha^{2p}).
\end{equation}
Here $A_\alpha=AE_\alpha$,
$B_\beta=BF_\beta$. In particular $[S,A]=[S,B]=0$.
Put $r=p/p_0$. For $r>1$, the first analytic family is
\[
 Z(z)=B^{rz}A^{rz}S,\qquad0\le\Re z\le1.
\]
On the operator boundary its norm is at most one. On the other
boundary,
$Z(1+it)=B^{irt}(B^rA^rS)A^{irt}$; this equality uses $[S,A]=0$.
Imaginary powers are contractions on their supports. Interpolating
at $1/r$ with exponent $2p_0$ on the right boundary gives
\[
 \|BAS\|_{2p}^{2p}\le\|B^rA^rS\|_{2p_0}^{2p_0}.
\]
For $r=1$ the assertion is equality. The projection resolutions give
$B^rA^rS=\sum_{\alpha,\beta}B_\beta^rA_\alpha^rS$.
Only the $J$ counted pairs contribute, since
$B_\beta^rA_\alpha^rS=B_\beta^r(F_\beta E_\alpha S)A_\alpha^r$.
At the fixed exponent $2p_0$, the triangle inequality bounds the last norm by
$J\max_{\alpha,\beta}\|B_\beta^rA_\alpha^rS\|_{2p_0}$.
For a pair of profiles use the second family
\[
 Z_{\alpha\beta}(z)=B_\beta^{pz}A_\alpha^{pz}S.
\]
At $z=it$ this equals
$B_\beta^{ipt}(F_\beta E_\alpha S)A_\alpha^{ipt}$, with norm at most
$w_{\alpha\beta}$. On the other boundary the outer imaginary powers
again disappear. Its squared Hilbert--Schmidt norm is at most
\[
 \operatorname{Tr}(S A_\alpha^pB_\beta^{2p}A_\alpha^pS)
 \le\operatorname{Tr}(B_\beta^{2p}A_\alpha^{2p}).
\]
Interpolation at $1/p_0$ therefore yields
\[
 \|B_\beta^rA_\alpha^rS\|_{2p_0}^{2p_0}
 \le w_{\alpha\beta}^{2(p_0-1)}
       \operatorname{Tr}(B_\beta^{2p}A_\alpha^{2p}).
\]
Complex powers remain zero on zero singular spaces throughout.
The trace on $S\mathcal H$ contains at least $L_\lambda$ identical
copies of the target block. Positivity of the singular power gives
$L_\lambda\operatorname{Tr}_\lambda|B_n|^{2p}
\le\|BAS\|_{2p}^{2p}$. These inequalities prove
\eqref{rec:hybrid-two-strip}. Both counting powers depend only on
$p_0$. In the tuple application $S$ is a central $S_n$-isotype;
in the signed-color application it is a global color isotype.
Both commute with the full permutation action and hence reduce
$A,B$ as required.

The finite starting range can be made entirely explicit.  Put
$n_0=2^{d_0}$ and let $\mathcal E_d$ be the hypercube edge
transpositions.  The identity and each member of $\mathcal E_d$ have
a specified sweep realization of probability $\varepsilon_n=2^{-nd/2}$.
For a unit vector $v$ in a nontrivial irreducible, the variance identity
for two independent sweep permutations gives
\[
 1-\|B_nv\|^2
 =\frac12\mathbb E_{g,h}\|\rho(g)v-\rho(h)v\|^2
 \ge\varepsilon_n^2\sum_{e\in\mathcal E_d}
       \|(\rho(e)-I)v\|^2.
\]
Every permutation is a word of at most $2n^2$ edge transpositions:
a vertex transposition uses a path of length at most $d$ and its return,
and at most $n-1$ vertex transpositions suffice.  Telescoping along such
a word and using Cauchy--Schwarz, then averaging uniformly over $S_n$,
gives
\[
 2=\mathbb E_{g\in S_n}\|\rho(g)v-v\|^2
 \le(2n^2)^2\sum_{e\in\mathcal E_d}\|(\rho(e)-I)v\|^2.
\]
The equality uses the absence of invariant vectors in a nontrivial
irreducible.  Thus the exact squared-norm gap is
\[
 \|B_n\|_{\mathrm{op}}^2\le1-\frac{2^{-nd}}{2n^4}.
\]
It is therefore enough to take, after the absolute threshold $d_0$
has been fixed,
\begin{equation}\label{rec:hybrid-explicit-base}
 p\ge2^{n_0d_0+1}n_0^5(\Lambda+2)\log n_0.
\end{equation}
Indeed, with $\Gamma_0=2^{-n_0d_0}/(2n_0^4)$, every nontrivial
starting block has squared norm at most $e^{-\Gamma_0}$.  The
unnormalized trace has $D_\lambda$ eigenvalues, and
$D_\lambda\le n!\le n_0^{n_0}$, so
\[
 \log T_n^\lambda(p)\le2n_0\log n_0-p\Gamma_0
 \le-\Lambda n_0\log n_0.
\]
The level-budget exponent is at least $-\Lambda n_0\log n_0$ on
these sizes, and the diagram budget is nonnegative; zero blocks are
automatic.  This deliberately crude choice therefore proves both
starting targets and makes the order of choices noncircular.

\subsection{Small levels by path cancellation}

For all sufficiently large $d$ and $1\le k\le n^{3/5}$ we claim
\begin{equation}\label{rec:hybrid-sparse}
 \|B_n|_{V_\lambda}\|_{\op}\le e^{-c_s k\log n}
\end{equation}
for an absolute $c_s>0$. We prove this on the injective $k$-tuple
module. Its $\lambda$-isotype has multiplicity
$D_{\bar\lambda}$, where $\bar\lambda$ is the diagram below the
first row. Zero extension to all $k$-tuples, with uniform product
measure, multiplies every squared norm by the same factor
$(n)_k/n^k$. An isotype vector has zero average in each separate
coordinate: otherwise its averaging would give a nonzero copy of
$\lambda$ in the $(k-1)$-tuple module, contrary to branching.

Let $x,y$ be independent uniform $k$-tuples. For $I\subseteq[k]$,
let $g_I(x,y)$ be $n^{|I|}$ times the endpoint transition
probability for the labels in $I$, extended by zero when either
restricted tuple is noninjective and then lifted to all coordinates.
Set $g_\varnothing=1$. On injective $x,y$, define a potential contact
between two labels as follows. If $t$ is their last differing input
bit in sweep order, their output prefixes of length $t-1$ must agree.
Their prescribed paths then enter opposite sides of the same switch
at time $t$. The prescribed routes are incompatible if their output
bits also agree at time $t$. Otherwise the shared switch saves one
of the $|I|d$ independent coin requirements. It follows that
$g_I=0$ on incompatible routes and
$g_I=2^{e_I}$ on compatible routes, where $e_I$ counts potential
contacts within $I$. Indeed, the input suffix and output prefix
specify the position before every layer; a collision of prescribed
positions would already give an incompatible potential contact.

In a bilinear form between the zero-extended isotype vectors, we
may replace $g_{[k]}$ by
$\sum_{I\subseteq[k]}(-1)^{k-|I|}g_I$. Each omitted-coordinate
term vanishes by the preceding zero-average property. On injective
$x,y$ this centered kernel vanishes unless the potential-contact
graph has no isolated vertex: toggle any isolated label in the sum.
Its squared Hilbert--Schmidt norm is consequently at most
\begin{equation}\label{s:hybrid-centered-kernel-square}
 4^k\max_{I\subseteq[k]}
 \mathbb E_{x,y}\!\left[
  {\bf1}_{\{x,y\ {\rm injective}\}}
  {\bf1}_{\{\text{every label has a potential contact}\}}g_I^2
 \right].
\end{equation}
All expectations in this display use uniform product measure.

Use one factor of $g_I$ as a change of measure. Namely, sample the
$q=|I|$ input sites uniformly without replacement and move them
through an independent complete switch field $\omega$; labels
outside $I$ retain independent uniform inputs and outputs. The
probability that the original independent inputs in $I$ were
injective is at most one, so it can be discarded in an upper bound.
The remaining factor is $2^{e_I}$, now counting encounters of the
$q$ sampled actual paths. We may also discard injectivity conditions
on outside labels. When their inputs coincide we define no potential
edge between them; this convention agrees with the original event
whenever all inputs are injective.

Fix $\omega$. Its encounter graph $\mathcal G_\omega$ has the $n$
initial sites as vertices; two vertices are adjacent when their
actual paths share a switch. A path has at most one such contact per
layer, so the degree is at most $d$. Two paths meet at most once:
after a meeting they differ in the edited bit, which is never edited
again. For any $s$ vertices there are at most $s\log_2s/2$ induced
edges. To verify this, merge the input suffix classes in a binary
tree. At a merge with selected child counts $h,s-h$, the new contacts
form a matching, of size at most $\min(h,s-h)$. The inequality
\[
 2\min(h,s-h)\le
 s\log_2s-h\log_2h-(s-h)\log_2(s-h)
\]
telescopes over the tree. Thus a sampled component of size $s$ has
weight at most $s^{s/2}$.

For the uniform injective sample of size $q\le k$, let $L$ be the
number of sampled vertices in nonsingleton induced components, let
$J$ be the number of these components, and let $E_I$ be the number
of sampled induced edges. We prove the precise weighted estimate
\begin{equation}\label{rec:hybrid-component-count}
 \mathbb E\!\left[{\bf1}_{\{L=l,J=j\}}2^{E_I}
                      \,\middle|\,\omega\right]
 \le e^{Ck\log(d+1)}\frac{k^{3l/2-2j}}{n^{l-j}}.
\end{equation}
The case $l=j=0$ has right-hand side at least one. For $j>0$,
enumerate ordered disjoint connected vertex sets of sizes
$s_1,\ldots,s_j\ge2$, with sum $l$. There are at most $2^l$ size
lists. A connected set of size $s$ can be encoded by a rooted plane
spanning tree, a root vertex, and a neighbor index for each tree
edge. There are at most $4^{s-1}$ plane tree shapes, so at most
$n(4d)^{s-1}$ such encodings. An encoding may be redundant; it is
used only for an upper bound.

Assign distinct sample labels to the chosen vertices in at most
$k^l$ ways. Prescribed distinct placements of those labels have
probability $1/(n)_l$. On the event $L=l,J=j$, the actual components
give at least $j!$ ordered witnesses. Each such witness, weighted by
$\prod_h s_h^{s_h/2}$, bounds $2^{E_I}$. Hence the expectation is
at most
\[
 \frac{2^l n^j(4d)^{l-j}k^l}{j!(n)_l}
       \max_{s_1+\cdots+s_j=l}\prod_hs_h^{s_h/2}.
\]
For $2\le s\le k$,
$s^{s/2}=s\,s^{(s-2)/2}\le e^s k^{(s-2)/2}$, so the product
is at most $e^l k^{(l-2j)/2}$. Moreover
$(n)_l\ge(n/e)^l$ and
$1/j!\le e^{Ck}k^{-j}$, the latter following from
$j!\ge(j/e)^j$ and $j\log(k/j)\le k/e$.
These bounds prove \eqref{rec:hybrid-component-count}, uniformly
in the fixed switch field.

Put $z=\log k/\log n\le3/5$. The logarithm of the ratio in
\eqref{rec:hybrid-component-count}, divided by $\log n$, is
\[
 (3z/2-1)l+(1-2z)j\le-l/10.
\]
For $z\ge1/2$ both coefficients have the required sign; for
$z<1/2$, use $j\le l/2$ to obtain the stronger bound
$-(1-z)l/2$. Thus $l\ge k/4$ supplies a saving of order
$k\log n$.

If $l<k/4$ and $q\ge k/2$, expose the outside endpoints before
sampling the $q$ base paths. Each prescribed outside path has at
most $d$ possible base-path neighbors: at every layer exactly one
base path occupies the opposite pre-switch site. Coverage therefore
requires at least $h=\lceil k/4\rceil$ further sampled labels,
outside the nonsingleton components, to belong to a set of at most
$kd$ initial sites. In the preceding witness count, choose these
labels in at most $2^k$ ways and prescribe their sites in at most
$(kd)^h$ ways, using $1/(n)_{l+h}$ for the joint distinct placements.
Since $l+h\le q\le k$, this inserts the factor
$e^{O(k)}(kd/n)^h$ in the same weighted bound.

If $q<k/2$, condition on all sampled trajectories. Let $o=k-q$.
To cover the outside vertices choose a rooted forest whose components
either end at a sampled label or have one free outside root.
Every component of the second kind contains at least two outside
vertices, so there are at most $o/2$ free roots and at least $o/2$
forest edges. Root choices cost at most $2^o$ and each other vertex
has at most $k$ parent choices. Expose each parent before its child.
For a fixed contact time, the child's required input suffix,
opposite input bit, and output prefix jointly have probability
$1/n$; summing over times gives $d/n$. Thus, uniformly in the sampled
trajectories, coverage costs at most
$2^o(kd/n)^{o/2}$ for all large $d$.

In each case the additional power of $kd/n$ or the bound
$n^{-l/10}$ saves a fixed multiple of $k\log n$.
The costs $O(k\log(d+1))$, the at most $(k+1)^2$ choices of $l,j$,
and the $4^k$ in \eqref{s:hybrid-centered-kernel-square} are smaller
than that saving when $d$ is large. The resulting Hilbert--Schmidt
bound for the centered kernel proves \eqref{rec:hybrid-sparse},
after decreasing $c_s$. When $k=1$ the centered kernel is zero.
Finally $D_\lambda\le\binom nkD_{\bar\lambda}\le n^k$, so
$T_n^\lambda(p)\le n^{2k}e^{-2pc_s k\log n}$.
Choosing a fixed $p$ sufficiently large in terms of $\Lambda,c_s$
proves the level budget in this range; it also proves the nonnegative
diagram budget there.

\subsection{A tuple frame that includes full rows}

The sparse estimate settles the small levels. For larger levels we need
both the angle and the trace endpoint in~\eqref{rec:hybrid-two-strip}.
We first construct a frame with total norm cost $e^{Ck}$; the entropy
saving will come from centering the resulting pieces.

For definiteness take the first bit batch to have $\lfloor d/2\rfloor$
bits. Thus the grid has $a=2^{\lceil d/2\rceil}$ rows and
$b=2^{\lfloor d/2\rfloor}$ columns; put $m=b$.
For larger $k$, work on the injective $k$-tuple module. Here $E_\alpha$
and $F_\beta$ are products of the local row and column Specht-isotypic
projections, with full shape profiles $\alpha=(\mu_i)_i$ and
$\beta=(\nu_j)_j$, where $\mu_i\vdash b$ and $\nu_j\vdash a$.
Let the local levels be $s_i=b-(\mu_i)_1$ and $t_j=a-(\nu_j)_1$,
with totals $J_s,J_t\le k$, and put
$C_1=\sum_i s_i\log(b/s_i)+\sum_jt_j\log(a/t_j)$.
The required angle estimate is
\begin{equation}\label{rec:hybrid-tuple-angle}
 \log w\le Ck-\tfrac14\{k\log(n/k)-C_1\}_+.
\end{equation}
We give the full frame construction, including completely occupied
rows.  Fix the tuple's row assignment.  In one row let $b$ be the
number of sites, $u$ the number of tuple labels, and
$s=b-\mu_1$ the local level.  Branching requires
$0\le s\le u\le b$.  Identify its observed labels with $[u]$ and
embed its function isometrically in $L^2(S_b)$ by completing to
$b$ labels and ignoring those outside $[u]$.  All measures here are
uniform probability measures.

For $Z\subseteq[b]$, let $P_Z$ be conditional averaging given the
values at labels in $Z$.  This is a right-subgroup average, commuting
with the left action on values.  The sum over $|Z|=s$ is right central.
In the regular representation, the whole left $\mu$-isotypic space
is $V_\mu\otimes V_\mu^*$, with right action on the paired carrier.
The invariant dimension for a subgroup permuting the $b-s$
complementary labels is $D_{\bar\mu}$ by branching, where
$\bar\mu$ is the lower diagram.  Taking the trace on that right
carrier therefore proves
\begin{equation}\label{rec:hybrid-frame-scalar}
 \sum_{|Z|=s}P_Z=\kappa I,\qquad
 \kappa=\binom bs\frac{D_{\bar\mu}}{D_\mu}\ge1
\end{equation}
on the entire left isotypic space, including its multiplicity.
The inequality follows by choosing the entries below the first row
of a standard tableau.  Each $P_Zf$ remains in this left type and
is harmonic in its $s$ selected values: a function of fewer selected
labels cannot contain a type of level $s$.

Since $P_{[u]}f=f$, recover the original function by applying
$P_{[u]}$ to \eqref{rec:hybrid-frame-scalar}.  Put
$l=|Z\setminus[u]|$ and $v=b-u$.  The only possible indices satisfy
\begin{equation}\label{rec:hybrid-frame-feasible}
 0\le l\le\min(s,v),\qquad |Z\cap[u]|=s-l.
\end{equation}
For such $l$, conditional averaging of the unobserved values in
$P_Zf$ is precisely
\begin{equation}\label{rec:hybrid-frame-reassignment}
 P_{[u]}P_Zf
 =\frac{(-1)^l}{(v)_l}
 \sum_{\varphi:Z\setminus[u]\hookrightarrow[u]\setminus Z}
          (P_Zf)\big|_{x_z=x_{\varphi(z)},\ z\in Z\setminus[u]}.
\end{equation}
To prove this identity, average the $l$ unobserved arguments
injectively over the values not used by the observed $u$ labels.
Starting with distinctness among the selected arguments, use
inclusion--exclusion to forbid also the values at $[u]\setminus Z$.
Every term leaving an argument unassigned to one of these values
vanishes by harmonicity.  The surviving assignments are exactly the
injections $\varphi$ in the display, each with sign $(-1)^l$.
Their number is $(u-s+l)_l\le(u)_l$.  Every reassigned function
has the same norm as $P_Zf$, since its $s$ distinct observed values
have the same uniform injective marginal distribution.

Use $(v)_0=(u)_0=1$, also when $v=0$.  For $l>0$ in
\eqref{rec:hybrid-frame-feasible}, $v\ge l$, so the denominator
is strictly positive.  In particular, for a full row $u=b$ only
$l=0$ occurs; no quotient with a zero falling factorial is evaluated.
By Cauchy--Schwarz and
$\sum_{|Z|=s}\|P_Zf\|^2=\kappa\|f\|^2$, the sum of the norms
of the reassigned pieces, after division by $\kappa$, is bounded by
$C_{b,u,s}\|f\|$, where
\begin{equation}\label{rec:hybrid-frame-cost}
 C_{b,u,s}=\kappa^{-1/2}
 \left[\sum_{l=0}^{\min(s,v)}
      \binom u{s-l}\binom vl
      \left(\frac{(u)_l}{(v)_l}\right)^2\right]^{1/2}
 \le e^{Cu}.
\end{equation}
An empty row has $u=s=0$ and cost one.  For a full row, the expression
is $\kappa^{-1/2}\binom bs^{1/2}\le2^{u/2}$.
More generally, if $v<2u$, then for every feasible $l$
$ (u)_l/(v)_l=\binom ul/\binom vl\le2^u$.
Vandermonde's identity bounds the sum of the binomial factors by
$\binom bs\le2^b\le2^{3u}$, proving the stated cost in this
dense case.  If $v\ge2u$, then $l\le u\le v/2$ and, for $l>0$,
\[
 \binom vl\left(\frac{(u)_l}{(v)_l}\right)^2
       \le\left(\frac{4e u^2}{vl}\right)^l.
\]
Using $\binom u{s-l}\le2^u$ and
$\sum_{l\ge0}(A/l)^l\le e^A$ with the $l=0$ term defined as
one (termwise, $l!\le l^l$), where $A=4eu^2/v\le2eu$,
proves \eqref{rec:hybrid-frame-cost} in the sparse case as well.

The construction also applies to entangled row vectors.  With the
fixed row assignment denoted by $\mathbf u$, take
$P_{\mathbf Z}=\prod_iP_{Z_i}$ on the completed-label product space.
The product of the central sums acts as $\prod_i\kappa_i$ times
identity on the full product isotypic space.  Consequently
\[
 f_{\mathbf u}=(\prod_i\kappa_i)^{-1}
       \sum_{\mathbf Z}(\prod_iP_{[u_i]})P_{\mathbf Z}f_{\mathbf u},
 \qquad
 \sum_{\mathbf Z}\|P_{\mathbf Z}f_{\mathbf u}\|^2
      =(\prod_i\kappa_i)\|f_{\mathbf u}\|^2.
\]
Conditional averaging and harmonicity hold with all other row
variables fixed.  Weighted Cauchy--Schwarz in this identity therefore
multiplies the costs in \eqref{rec:hybrid-frame-cost}, without any
assumption that $f_{\mathbf u}$ is a tensor product.  Since
$\sum_i u_i=k$, the total cost is $e^{Ck}$.

Group the reassigned pieces by their global exceptional subset
$X\subseteq[k]$.  In the independent-coordinate extension they have
the form $H_X(\mathbf u,\mathbf v_X)\mathbf1_{\rm distinct}$,
with $|X|=J_s$ and at most $2^k$ subsets. In each supported row
pattern, exactly $s_i$ labels of $X$ lie in row $i$.
Set the raw function to zero on unsupported
row patterns and on collisions among its selected columns.
Conditional on a row pattern, the selected injective columns have
the same marginals as under full row injectivity.  Moreover the
probability of full row injectivity is
$\prod_i(b)_{u_i}/b^{u_i}\ge e^{-k}$.
Thus the raw independent-coordinate norms still cost at most
$e^{Ck}$.  A column-isotype vector has the analogous pieces
$G_Y(\mathbf u_Y,\mathbf v)\mathbf1_{\rm distinct}$, where
$|Y|=J_t$ and exactly $t_j$ labels of $Y$ lie in column $j$.

\subsection{Conditional collisions and localization}

The frame has paid only $e^{Ck}$ and has not yet produced the saving
in~\eqref{rec:hybrid-tuple-angle}. We now use centering to force
collisions among omitted labels, and localize the selected labels
using their row and column counts.

The vectors in the global level-$k$ type are harmonic after zero
extension.  Insert the centering projection
$\Pi=\prod_{\ell=1}^k(I-M_\ell)$ into their inner product, where
$M_\ell$ averages the independent cell of label $\ell$.
For labels in $X\cup Y$, expand this product and let
$I\subseteq X\cup Y$ be the subset on which the identity factor
is chosen. Write the cells supplied to the two functions as
$(U_\ell^H,V_\ell^H)$ and $(U_\ell^G,V_\ell^G)$.
For $\ell\in I$ these are the same uniform cell. For
$\ell\in(X\cup Y)\setminus I$ they are independent uniform cells.
Different labels are independent. In particular a label in
$(X\cap Y)\setminus I$ supplies two independent selected cells;
it is not a shared selected cell.

For each missing label in $Z=[k]\setminus(X\cup Y)$, couple its
two choices using main row and column variables $(U_\ell,V_\ell)$
and independent alternates $(U'_\ell,V'_\ell)$.  The identity choice
uses $(U_\ell,V_\ell)$ on both sides; the averaging choice uses
$(U_\ell,V'_\ell)$ in $H_X$ and $(U'_\ell,V_\ell)$ in $G_Y$.
The raw functions always see the same needed row $U_\ell$ and column
$V_\ell$.  Hence toggling a missing label changes only the two
distinctness indicators.  Their alternating sum cancels unless that
label has a possible same-cell collision in some choice of side and
cell version.
Put $M=|Z|$ and condition first on all nonmissing variables.
The raw functions have separated remaining variables $H(U),G(V)$.
Whenever the alternating sum is nonzero, the graph of possible
collisions has a forest covering the missing vertices. In every
component meeting nonmissing vertices, grow a forest from all those
vertices as separate roots; then each tree has exactly one nonmissing
root. Span each wholly missing component by one tree, which has at
least two vertices and needs one free root. Thus every child exposed
along a forest edge is missing, and there are at least $M/2$ edges.
Choosing parents, sides, and cell versions costs
$e^{O(k)}k^{\text{number of edges}}$. Exposing parent before child,
each specified child cell is independent uniform, so every edge
costs $1/n$. Uniformly under the conditioning, the coverage probability
is at most $e^{O(k)}(k/n)^{M/2}$.

Let $K(U,V)$ be the alternating sum after averaging alternate
coordinates. Jensen's inequality and its $2^M$ summands give
$\mathbb E|K|^2\le4^M\Pr(\text{coverage})$.
It vanishes unless the cells in $O=X\cap Y\cap I$ are distinct.
Writing $A_1=(\mathbb E_U|H|^2)^{1/2}$ and
$B_1=(\mathbb E_V|G|^2)^{1/2}$, the conditional Cauchy--Schwarz
bound is therefore
\begin{equation}\label{s:conditional-collision-weight}
 e^{O(k)}(k/n)^{M/4}{\bf1}_{\{O\text{ distinct}\}}A_1B_1.
\end{equation}
Here the factorization
$\mathbb E_{U,V}|H(U)G(V)|^2=A_1^2B_1^2$ is essential:
an unconditional probability estimate alone would not control the
weighted functions.

The dependence of these two weights will also matter in the next step.
Put $N=X\cup Y$. After averaging the missing rows, $A_1$ is measurable
with respect to
\[
 \mathcal F_H=\sigma\big((U_\ell^H)_{\ell\in N},
                         (V_\ell^H)_{\ell\in X}\big).
\]
Similarly, after averaging the missing columns, $B_1$ is measurable
with respect to
\[
 \mathcal F_G=\sigma\big((V_\ell^G)_{\ell\in N},
                         (U_\ell^G)_{\ell\in Y}\big).
\]
Only the cells in $O=X\cap Y\cap I$ are selected on both sides
with a shared value. Given $\mathcal F_H$, the columns
$(V_\ell^G)_{\ell\in Y\setminus O}$ are still independent and
uniform; given $\mathcal F_G$, the rows
$(U_\ell^H)_{\ell\in X\setminus O}$ have the analogous property.
For a label selected on both sides but outside $I$, this follows from
the independent versions just specified.

Stratify the nonmissing variables by the row and column counts
$p_i,q_j$ of $O$, whose size is $o$. There are at most
$(k+1)^{a+b}=e^{O(k)}$ strata because $k\ge n^{3/5}$ and the grid is
balanced. On a realizable stratum $p_i\le s_i,q_j\le t_j$.
For $o>0$, the entropy of the uniform measure on its distinct cells
is $\log o$, bounded by the sum of its marginal entropies. Hence
\[
 \sum_i p_i\log(b/p_i)+\sum_jq_j\log(a/q_j)
 \ge o\log(n/o).
\]
Using $\sum_i p_i\log(s_i/p_i)\le J_s/e$ and its column analogue,
with zero terms interpreted continuously, gives
\begin{equation}\label{s:overlap-marginal-entropy}
 \sum_i p_i\log(b/s_i)+\sum_jq_j\log(a/t_j)
 \ge o\log(n/k)-2k/e.
\end{equation}
This is also true when $o=0$.

We make the localization step explicit. In the average of $A_1B_1$
on a stratum, apply Cauchy--Schwarz, keeping the stratum and the
$Y$-column profile in the $A_1^2$ factor and the stratum and the
$X$-row profile in the $B_1^2$ factor. Given $\mathcal F_H$,
the cells of $O$ are known and the columns seen by $G$ on
$Y\setminus O$ remain independent uniform columns. The probability
of the remaining column counts $r_j=t_j-q_j$, with $N_Y=|Y|-o$, is
multinomial. Its logarithm is at most
\begin{align*}
 \log\frac{N_Y!}{\prod_jr_j!}-N_Y\log b
 &\le-N_Y\log(n/k)+\sum_jr_j\log(a/t_j)+O(k).
\end{align*}
Indeed $N_Y!\le k^{N_Y}$, $r!\ge(r/e)^r$, and
$\sum_jr_j\log(t_j/r_j)\le J_t/e$. For the other factor the same
argument gives the probability bound
\[
 \exp\left\{-(|X|-o)\log(n/k)
       +\sum_i(s_i-p_i)\log(b/s_i)+O(k)\right\}.
\]
The unweighted second moments are the squared raw norms. Taking
square roots of the two probabilities and using
\eqref{s:overlap-marginal-entropy} bounds the localization factor by
\[
 \|H\|_2\|G\|_2
 \exp\left\{O(k)-\tfrac12[(|X|+|Y|-o)\log(n/k)-C_1]\right\}.
\]
Also $|X|+|Y|-o\ge|X\cup Y|=k-M$, and the trivial
Cauchy--Schwarz bound is available. We thus obtain the positive-part
saving $\tfrac12((k-M)\log(n/k)-C_1)_+$.
Finally, for $L=\log(n/k)\ge0$,
\[
 ML/4+((k-M)L-C_1)_+/2\ge(kL-C_1)_+/4.
\]
Combining this with \eqref{s:conditional-collision-weight}, and summing
the strata, centering choices, and frame pieces, proves
\eqref{rec:hybrid-tuple-angle}.

\subsection{The tuple trace endpoint}

The angle estimate~\eqref{rec:hybrid-tuple-angle} is now complete.
It remains to bound the positive trace endpoint, including the global
copy count and all row and column assignment factors.
We now estimate the trace endpoint in
\eqref{rec:hybrid-two-strip}. Write $n=ab$, with $a$ rows of size
$b$ and $b$ columns of size $a$. Put
$X=A_\alpha^{2p}$ and $Y=B_\beta^{2p}$ on the injective tuple
module. In a diagonal entry of $YX$, an intermediate tuple must
have both the starting row assignment, because $X$ preserves rows,
and the starting column assignment, because $Y$ preserves columns.
It is therefore the starting tuple itself. Consequently
\[
 \Tr(YX)=\sum_{z\ {\rm injective}}Y(z,z)X(z,z),
\]
and both diagonals are nonnegative.

For a tuple with column counts $v_j$, consider the local positive
group-algebra element of $Y$ in column $j$. Its identity coefficient
is $T_a^{\nu_j}(p)/a!$ and the modulus of every coefficient is at
most this value, by positivity in the regular representation.
The diagonal entry in the local tuple action sums coefficients on
the pointwise stabilizer of its $v_j$ labels, of size $(a-v_j)!$.
It is thus at most $T_a^{\nu_j}(p)/(a)_{v_j}$. Equivalently this
is positive coefficient restriction to that stabilizer.
Since $\sum_jv_j=k$ and $(a)_v\ge(a/e)^v$,
\begin{equation}\label{s:hybrid-tuple-column-diagonal}
 Y(z,z)\le(e/a)^k\prod_jT_a^{\nu_j}(p).
\end{equation}

Fix the complete row assignment of the labeled tuple, with counts
$u_i$. The local row tuple module is
$\operatorname{Ind}_{S_{b-u_i}}^{S_b}{\bf1}$, in which the
multiplicity of $\mu_i=(b-s_i,\bar\mu_i)$ is
$f^{\mu_i/(b-u_i)}$. It is zero if $u_i<s_i$. Otherwise separating
the entries below the first row gives
\[
 f^{\mu_i/(b-u_i)}\le\binom{u_i}{s_i}D_{\bar\mu_i}.
\]
The hook formula also gives
\begin{equation}\label{s:hybrid-first-row-dimension}
 D_{\mu_i}\ge e^{-s_i}\binom b{s_i}D_{\bar\mu_i}.
\end{equation}
Indeed the hooks below the first row are exactly those of
$\bar\mu_i$. If $r=b-s_i$, the first-row hook product divided by
$r!$ is
$\prod_{h=1}^r(1+(\bar\mu_i)'_h/(r-h+1))$; its logarithm is at
most $\sum_h(\bar\mu_i)'_h=s_i$.

The trace of the local row operator is its irreducible trace times
this multiplicity. Hence, for the fixed row assignment, its total
row trace equals
\[
 \prod_i T_b^{\mu_i}(p)
       \frac{f^{\mu_i/(b-u_i)}}{D_{\mu_i}}.
\]
Using \eqref{s:hybrid-first-row-dimension},
$\binom b{s}\ge(b/s)^s$, and
$\binom u{s}\le(eu/s)^s$, the product of ratios is at most
\[
 \exp\left\{-C_H+O(k)\right\},\qquad
 C_H=\sum_i s_i\log(b/s_i).
\]
Here $\sum_i s_i\log(u_i/s_i)\le\sum_i u_i/e=k/e$ and
$\sum_i s_i\le k$; zero terms are interpreted continuously.

The number of row assignments with every $u_i\ge s_i$ is at most
\begin{equation}\label{s:hybrid-row-assignment-count}
 a^k\exp\{C_H-J_s\log(n/k)+O(k)\}.
\end{equation}
For a uniform independent row assignment, choose disjoint witness
sets of $s_i$ labels required to hit row $i$. There are at most
$k^{J_s}/\prod_i s_i!$ such choices, each of probability
$a^{-J_s}$. The union bound and $s!\ge(s/e)^s$ give
\[
 \log\Pr(u_i\ge s_i\ \forall i)
 \le J_s\log(k/a)-\sum_i s_i\log s_i+O(k)
 = C_H-J_s\log(n/k)+O(k),
\]
which proves \eqref{s:hybrid-row-assignment-count}.

Combine the trace per assignment, the count of assignments, and
\eqref{s:hybrid-tuple-column-diagonal}. The $a^k$ cancels and gives
\[
 \Tr(YX)\le
 \left(\prod_iT_b^{\mu_i}(p)\prod_jT_a^{\nu_j}(p)\right)
       e^{-J_s\log(n/k)+O(k)}.
\]
Interchanging rows and columns gives the same inequality with $J_t$.
The global tuple multiplicity is $L_\lambda=D_{\bar\lambda}$,
and $D_\lambda/L_\lambda\le\binom nk\le(en/k)^k$.
Thus the full normalized endpoint satisfies
\begin{equation}\label{rec:hybrid-tuple-trace}
 \log\left(\frac{D_\lambda}{L_\lambda}
                      \Tr(YX)\right)
 \le\sum_i\log T_b^{\mu_i}(p)+\sum_j\log T_a^{\nu_j}(p)
       +(k-\max(J_s,J_t))\log(n/k)+O(k).
\end{equation}
Pairs with zero endpoint require no logarithmic estimate.
There are only $e^{O(k)}$ relevant profile pairs in this range.
Indeed the level lists have at most $(k+1)^{a+b}$ choices, and,
for fixed lists, the partition bound gives at most
$\exp(C\sum_i\sqrt{s_i}+C\sum_j\sqrt{t_j})$
tail choices. Cauchy--Schwarz bounds the exponent by
$O(\sqrt{ak}+\sqrt{bk})$. Since $a,b\asymp\sqrt n$ and
$k\ge n^{3/5}$, both this exponent and $(a+b)\log(k+1)$ are $O(k)$.

\subsection{Closing the level budget}

In both inductions, a trivial child has $T=1$ and zero level and
diagram budgets. If a child sweep block is zero, its profile has zero
trace endpoint and can be discarded. Every remaining child satisfies
the height condition by Lemma~\ref{lem:annihilation}, so the stated
induction hypotheses apply to all child terms retained below.

Insert the angle~\eqref{rec:hybrid-tuple-angle}, the endpoint
bound~\eqref{rec:hybrid-tuple-trace}, and the profile count into
\eqref{rec:hybrid-two-strip}. The remaining task is to compare the
sum of child level budgets with the parent budget.
Write
$\delta=\log(n/k)/\log n$, let
$\omega_i=(s_i/k)(\log b/\log n)$ for row factors and use the
corresponding column weights. For a positive local level $s_i$ in a
block of size $m_i$, put
$\delta_i=\log(m_i/s_i)/\log m_i$, where $m_i=b$ on rows and
$m_i=a$ on columns (with $s_i$ replaced by $t_i$ there).
At level zero set $\delta_i=0$; its weight is zero and the trivial
child moment is one. Put
\[
 \bar u=\max\{u_{\lfloor d/2\rfloor},u_{\lceil d/2\rceil}\},
 \qquad W_0=\sum\omega_i,
 \qquad \xi=(\delta-\sum\omega_i\delta_i)_+.
\]
Then
\[
 W_0\le\max(J_s,J_t)/k,\qquad
 \sum\omega_i\delta_i^2\ge(2-W_0)\delta^2-2\xi.
\]
For the second inequality, expand
$\sum_i\omega_i(\delta_i-\delta)^2\ge0$ and use
$\sum_i\omega_i\delta_i\ge\delta-\xi$. This gives
$\sum_i\omega_i\delta_i^2\ge(2-W_0)\delta^2-2\delta\xi$;
since $0\le\delta\le1$, it implies the displayed bound.
Subtracting the target level budget and dividing by $k\log n$
leaves at most
\begin{equation}\label{rec:hybrid-level-close}
 -(\bar u+\Lambda\delta^2-\delta)(1-W_0)
 +(2\Lambda-(p_0-1)/2)\xi
 -(u_d-\bar u)+O(p_0/\log n).
\end{equation}
The first two terms are nonpositive by
\eqref{rec:hybrid-constants}; the parameter increment, of order
$d^{-1/3}$ at a balanced split, pays the last error.  This proves
\eqref{rec:hybrid-level-budget}.

\subsection{Long chains and the diagram dimension}

It remains to treat $\delta<\epsilon$, so $k\ge n^{1-\epsilon}$.
For the two moment inequalities we may assume the parent block is
nonzero, since zero moments were handled above.  Throughout this
subsection the theorem's height condition is $\lambda'_1\le n/2$.
Put $T=n^{1-\gamma_d}$ and let $a_1,b_1$
be the numbers of full rows and full columns of $\lambda$ of length
at least $T$. Split the diagram into all boxes in its top $a_1$
rows, then all boxes in its first $b_1$ columns below those rows,
and the translated straight diagram $\theta$ remaining in the
lower right (Figure~\ref{fig:hybrid-chains}). The first two regions are disjoint chains, with lengths
$\eta_i=\lambda_i$ and $\zeta_j=(\lambda'_j-a_1)_+$, respectively.
Write
\[
 R=|\theta|,\qquad M=n-R,\qquad
 G=\sum_{l\ {\rm in}\ \eta,\zeta}l\log(M/l),\qquad
 D_0=1+a_1+b_1=O(n^{0.01}).
\]
The bound on $D_0$ follows from $a_1,b_1\le n/T=n^{\gamma_d}$
and $\gamma_d<0.01$. Set $G=0$ when $M=0$.

\begin{figure}[t]
\centering
\begin{tikzpicture}[x=0.38cm,y=0.38cm,font=\small]
  \foreach \row/\length in {1/9,2/8,3/6,4/5,5/5,6/3,7/2}{
    \foreach \col in {1,...,\length}{
      \ifnum\row<4
        \fill[blue!16] (\col-1,-\row+1) rectangle (\col,-\row);
      \else
        \ifnum\col<4
          \fill[green!18] (\col-1,-\row+1) rectangle (\col,-\row);
        \else
          \fill[orange!25] (\col-1,-\row+1) rectangle (\col,-\row);
        \fi
      \fi
      \draw[gray!65,line width=0.35pt]
        (\col-1,-\row+1) rectangle (\col,-\row);
    }
  }
  \draw[blue!65!black,line width=0.7pt]
    (-0.45,0)--(-0.65,0)--(-0.65,-3)--(-0.45,-3);
  \node[anchor=east,align=right] at (-0.9,-1.5)
    {top $a_1$ rows\\chain lengths $\eta_i$};
  \draw[green!45!black,line width=0.7pt]
    (0,-7.45)--(0,-7.65)--(3,-7.65)--(3,-7.45);
  \node[anchor=north,align=center] at (1.5,-7.95)
    {first $b_1$ columns below the cut\\chain lengths $\zeta_j$};
  \node[anchor=west,align=left] (tail) at (7,-4.5)
    {translated tail $\theta$\\$R$ boxes};
  \draw[-{Stealth[length=1.6mm]},gray!75] (tail.west)--(5,-4.5);
  \draw[dashed,gray!75] (0,-3)--(6,-3);
  \draw[dashed,gray!75] (3,-3)--(3,-7);
\end{tikzpicture}

\caption{The chain decomposition, shown schematically. The top $a_1$
rows are taken first; the column chains contain only the boxes below
those rows. The remaining lower-right region is a translated straight
diagram $\theta$. Thus the row chains, lower-column chains, and tail
are disjoint, and their sizes sum to $n$. The cut is schematic; the
asymptotic threshold is defined in the text.}
\label{fig:hybrid-chains}
\end{figure}

Every tail cell pays the cap $v_d\log n$, since both of its full
line lengths are below $T$. In a top row, replacing the maximum
of the full row and column lengths by the row length can change
the contribution only in the first $b_1$ columns, at most $a_1b_1$
cells. This has total cost $O(c_d a_1b_1\log(en))$.
Below the top rows, in a long column the maximum is its full
column height. Replacing that height $\zeta_j+a_1$ by its chain
length $\zeta_j$ changes the logarithmic sum by at most
\[
 \zeta_j\log(1+a_1/\zeta_j)\le a_1
\]
for each nonempty chain. Thus the chain contribution is
$c_d\sum_l l\log(n/l)$ up to the same stated error.
Finally $\sum_l l\log(n/l)=G+M\log(n/M)$ and
$M\log(n/M)\le R$. We have proved
\begin{equation}\label{rec:hybrid-arm-weight}
 \mathcal D_d(\lambda)
 =c_dG+v_dR\log n+O(c_d(R+D_0^2\log(en))).
\end{equation}

We also need a hook-product bound for an arbitrary diagram
$\mu\vdash L$, $L\le n$, using these same values of $a_1,b_1$.
Let $l$ be its top-row and lower-column chain lengths, and let
$\vartheta$ be its translated residual diagram of size $q$.
Then
\begin{equation}\label{rec:hybrid-arm-hook}
 D_\mu\ge\frac{L!D_\vartheta}{q!\prod_l l!}
             \exp\{-C(q+D_0^2\log^2(en))\}.
\end{equation}
Here is the hook comparison in detail. Tail hooks are exactly those
of $\vartheta$. On a top row, compare each hook with the baseline
number of cells remaining in that row, including the current cell.
Cells in the first $b_1$ columns cost at most
$O(a_1b_1\log(en))$ altogether. In the other columns, the extra
leg has at most $a_1$ boxes in the top region plus the height of the
corresponding tail column. The first addition costs
$O(a_1\log(en))$ per row by the harmonic sum; the two additions can
be separated using $\log(1+x+y)\le\log(1+x)+\log(1+y)$.

Let the tail have width $w$ and nonincreasing column heights
$h_1,\ldots,h_w$, of sum $q$. If $q>0$, every top row extends at
least through column $b_1+w$, so the tail addition for that row
is bounded by
\[
 \sum_{t=1}^w\log\left(1+\frac{h_t}{w-t+1}\right).
\]
When $w\le D_0\log(en)$, there are at most $a_1w$ affected top-row
cells, each costing at most $\log(en)$, for a total
$O(D_0^2\log^2(en))$. When $w>D_0\log(en)$, use
\[
 \sum_{t=1}^w\frac{h_t}{w-t+1}
 \le \frac{2q}{w}
       +\frac{2q}{w}\sum_{r=1}^{\lceil w/2\rceil}\frac1r
 =O\!\left(\frac q w\log(en)\right).
\]
For the first half of the sum the denominators are at least $w/2$;
for the second half monotonicity gives $h_t\le2q/w$.
After multiplication by the at most $a_1$ top rows, this is $O(q)$.

For the lower-column chains compare hooks with the number of
remaining cells in their column. The extra arm has at most $b_1$
boxes in the left region plus the appropriate tail row length.
The $b_1$ addition has total cost $O(b_1^2\log(en))$.
The tail addition has the identical bound with width and height
interchanged: every left column extends through the tail's rows,
and its row lengths are nonincreasing with sum $q$.
If $q=0$, both tail-addition terms are absent.
The resulting total logarithmic inflation over the product
$(q!/D_\vartheta)\prod_l l!$ is
$O(q+D_0^2\log^2(en))$, proving
\eqref{rec:hybrid-arm-hook}.

At the parent, every hook of $\theta$ is at most
$2n^{1-\gamma_d}$, because all its full row and column lengths
are below the threshold. Therefore
$D_\theta/R!\ge(2n^{1-\gamma_d})^{-R}$.
Also $(n)_R\ge(n/e)^R$ and factorial estimates give
$\log(M!/\prod_l l!)=G+O(D_0\log(en))$.
Applying \eqref{rec:hybrid-arm-hook} yields
\begin{equation}\label{rec:hybrid-arm-dimension}
 \log D_\lambda\ge
 G+\gamma_dR\log n-O(R+D_0^2\log^2(en)),
 \qquad G+R\log n\ge k/2 .
\end{equation}
For the second assertion, if $M>0$ and $l_{\max}$ is the largest
chain length, $-\log x\ge1-x$ gives
\[
 G\ge M-\frac{\sum_l l^2}{M}\ge M-l_{\max}.
\]
Every top chain has length at most $\lambda_1=n-k$, and every
lower-column chain has length at most $\lambda'_1\le n/2$.
Consequently
$G+R\ge n-l_{\max}\ge\min(k,n/2)\ge k/2$.
If $M=0$ the same inequality is immediate from $R=n$.
For sufficiently large $d$, $\log n\ge1$, proving the displayed
version.

\subsection{Color multiplicities and the marked trace}

The arm comparison has expressed the parent budget in terms of $G$
and $R$, and~\eqref{rec:hybrid-arm-dimension} supplies its dimension
scale. To apply~\eqref{rec:hybrid-two-strip}, we still need the local
copy counts, the positive trace with distinct markers, the child
diagram costs, and the color angle. We establish them in that order.

Let $\mathcal H$ be the space of configurations of $R$ distinctly
labeled even markers and signed colors in
$\mathbb C^{a_1}\oplus\mathbb C^{b_1}$. Every marker is used exactly
once. Thus $\mathcal H$ is the orthogonal sum, over the ordered
marker-position tuples, of the signed tensor spaces on the remaining
$M$ sites. The compact group $K=U(a_1)\times U(b_1)$ acts only on
colors; zero factors are omitted.  The selected global type is the
irreducible representation
\[
 \begin{gathered}
 \pi=\pi_{(\eta,\zeta)}:K\longrightarrow
 U(\mathcal U_{\eta,\zeta}),\\
 \mathcal U_{\eta,\zeta}
 =\mathbf S_\eta(\mathbb C^{a_1})
    \otimes\mathbf S_\zeta(\mathbb C^{b_1}),
 \end{gathered}
\]
where the carrier factor belonging to a zero group factor is also
omitted.  Its carrier action extends polynomially to the invertible
parity-preserving block matrices.  Take $S$ to be the full compact
$\pi$-isotypic projection on $\mathcal H$, including all multiplicity
copies. No additional projection is put on the marker labels.

Indeed, by fixing the even, odd, and marker site pattern and then
inducing, Schur--Weyl duality identifies this color isotype on the
permutation side with induction from the Specht factors
$V_\eta,V_{\zeta^{\mathsf t}}$, with the markers fixed pointwise.
Stacking $\zeta^{\mathsf t}$ below $\eta$ gives precisely the
unmarked hook subdiagram. Its Littlewood--Richardson coefficient
is at least one: the skew part below $\eta$ is a translated straight
diagram with its row-constant Littlewood--Richardson tableau.
Adding the $R$ ordered markers by branching then gives at least
$f^{\lambda/(\text{unmarked hook})}=D_\theta$ copies of $V_\lambda$.
The compact carrier dimension can only increase this multiplicity,
so in \eqref{rec:hybrid-two-strip} we may take
$L_\lambda=D_\theta$. Separating the entry sets in the chains and
the tail of a standard tableau gives
\begin{equation}\label{s:hybrid-global-marker-multiplicity}
 \frac{D_\lambda}{L_\lambda}
 \le \binom nR\frac{M!}{\prod_h l_h!}
 \le \binom nR e^G.
\end{equation}
The last inequality is the elementary multinomial entropy bound.

In each row include both the local permutation isotype $\mu_i$
and local color isotype $(\eta^{(i)},\zeta^{(i)})$ in $E_\alpha$; use
$(\nu_j,\eta^{[j]},\zeta^{[j]})$ in each column for $F_\beta$.
The superscripts $(i)$ and $[j]$ distinguish row and column
partitions; subscripts denote their parts. Local color projectors
commute with local permutations. They also commute with the global color projection
because they commute with every element of their local color
group, including the corresponding factor of every diagonal global
rotation. Thus these projections have all the commutations needed
in \eqref{rec:hybrid-two-strip}.
Their degrees fix the marker counts
\[
 x_i=b-|\eta^{(i)}|-|\zeta^{(i)}|,\qquad
 y_j=a-|\eta^{[j]}|-|\zeta^{[j]}|,\qquad
 \sum_i x_i=\sum_jy_j=R.
\]
Only nonnegative counts contribute. For $M_i=b-x_i$ define
\[
 G_i=\sum_{l\ {\rm in}\ \eta^{(i)},\zeta^{(i)}}l\log(M_i/l),
 \qquad G_H=\sum_iG_i,
\]
and define $G'_j,G_V$ analogously in columns. Empty entropies are zero.

We need the local multiplicity with one fixed assignment of marker
labels to a row, while their positions within that row still vary.
The unmarked permutation shapes $\rho\vdash M_i$ have coefficients
$c_{\eta^{(i)},(\zeta^{(i)})^{\mathsf t}}^\rho$. The full local multiplicity
of $\mu_i$ is therefore
\begin{equation}\label{s:hybrid-local-marker-multiplicity}
 \operatorname{mult}_i
 =(\dim\eta^{(i)})(\dim\zeta^{(i)})
   \sum_{\rho\subseteq\mu_i}
    c_{\eta^{(i)},(\zeta^{(i)})^{\mathsf t}}^\rho f^{\mu_i/\rho}.
\end{equation}
Here the two dimensions are compact-group carrier dimensions, and
the sum includes only shapes with nonzero coefficient.

For clarity, we record the size bounds used in this sum. Every such
$\rho$ contains $\eta^{(i)}$ and $(\zeta^{(i)})^{\mathsf t}$ and lies in the
$(a_1,b_1)$ hook. For the latter assertion use a Littlewood--Richardson
tableau of $\rho/(\zeta^{(i)})^{\mathsf t}$ with content $\eta^{(i)}$.
Beyond column $b_1$ its columns have at most $a_1$ boxes, by column
strictness. Let $l^0$ be the vector of its top-row and lower-column
chain lengths, padded by zeros. Its top counts dominate the parts
of $\eta^{(i)}$, while the deficit of each lower-column count from the
corresponding part of $\zeta^{(i)}$ is at most $a_1$. Since the two
vectors have the same total $M_i$, their total absolute discrepancy
is at most $2a_1b_1$. It follows that
\begin{equation}\label{s:hybrid-local-chain-entropy}
 \sum_h l_h^0\log(M_i/l_h^0)
       =G_i+O(D_0^2\log(en)).
\end{equation}
One may bound the entropy change per changed integer unit by
$O(\log(en))$, including the transition between zero and one.

The same tableau argument bounds the coefficient in
\eqref{s:hybrid-local-marker-multiplicity}. The top $a_1$ rows of
$\rho$ contain at least $|\eta^{(i)}|$ boxes, and at most $a_1b_1$
of those belong to the base shape $(\zeta^{(i)})^{\mathsf t}$.
Thus at most $a_1b_1$ boxes of the skew tableau lie below the top
region. In the top region, a weakly increasing row is specified by
its counts of the at most $a_1$ letters, giving at most
$(n+1)^{a_1^2}$ choices. Below it there are at most
$\max(1,a_1)^{a_1b_1}$ choices. The number of possible hook shapes
$\rho$ is at most $(n+1)^{a_1+b_1}$, and the compact carrier
dimensions are $e^{O(D_0^2\log(en))}$ by the Weyl dimension formula.
All these factors therefore fit that exponential bound. In particular,
\eqref{rec:hybrid-arm-hook} and
\eqref{s:hybrid-local-chain-entropy} give
\begin{equation}\label{s:hybrid-unmarked-dimension}
 \log D_\rho\ge G_i-O(D_0^2\log^2(en)).
\end{equation}

For each contributing $\rho\subseteq\mu_i$, let $l_{ih}$ be the
full chain lengths of $\mu_i$, and let $l^0_{ih}$ be those of
$\rho$. Write
\[
 e_{ih}=l_{ih}-l^0_{ih}\ge0,\qquad
 K_i(\rho)=\sum_h\log\binom{l_{ih}}{e_{ih}}.
\]
If the tail of $\mu_i$ has shape $\theta_i$ and size $q_i$, then
$q_i+\sum_he_{ih}=x_i$. Separating the entries in its tail and
chain extensions yields
\[
 f^{\mu_i/\rho}\le
       \frac{x_i!D_{\theta_i}}{q_i!\prod_he_{ih}!}.
\]
Dividing by \eqref{rec:hybrid-arm-hook}, with $L=b$, leaves
$x_i!\prod_hl_{ih}!/(b!\prod_he_{ih}!)$ times
$e^{O(x_i+D_0^2\log^2(en))}$. Factor out the chain binomials and
use \eqref{s:hybrid-local-chain-entropy} to obtain
\begin{equation}\label{s:hybrid-local-ratio}
 \frac{\operatorname{mult}_i}{D_{\mu_i}}
 \le\binom b{x_i}^{-1}
       \exp\{K_i-G_i+O(x_i+D_0^2\log^2(en))\},
\end{equation}
where, once and for all, choose a contributing $\rho_i$ maximizing
$K_i(\rho)$ and set $K_i=K_i(\rho_i)$.
The identity behind the entropy term is
$\prod_h(l^0_{ih})!/M_i!$ after extracting the inverse binomial;
factorial errors and the sum in
\eqref{s:hybrid-local-marker-multiplicity} have the stated size.
Zero local spaces can be discarded.

\paragraph{The marked trace endpoint.}
In a diagonal block of $B_\beta^{2p}A_\alpha^{2p}$, each marker's
intermediate position must equal its starting position: a horizontal
move and a vertical move cannot undo each other's coordinate changes.
Fix one such marker-position tuple. In a column, restricting the
positive group-algebra coefficients to its pointwise marker
stabilizer $S_{a-y_j}$ gives a positive element with regular trace
$T_a^{\nu_j}(p)/(a)_{y_j}$. On a stabilizer irreducible $\rho$,
its operator norm is at most this regular trace divided by $D_\rho$.
On the fixed local color type the occurring $\rho$ satisfy the
column version of \eqref{s:hybrid-unmarked-dimension}.
Consequently the vertical diagonal block has norm at most
\[
 (e/a)^R
 \exp\{-G_V+O(bD_0^2\log^2(en))\}
 \prod_jT_a^{\nu_j}(p).
\]
This holds uniformly over the marker positions; blocks with
incompatible counts are zero. Fixed graded reordering into column
order makes the product over columns literal, without changing
the norm or the marker fibers.

Both diagonal compressions are positive. After applying the vertical
norm bound we may sum the horizontal diagonal traces, which gives
the full horizontal trace. There are $R!/\prod_i x_i!$ ways to
assign the distinct marker labels to rows, and hence
\[
 \Tr_{\mathcal H}A_\alpha^{2p}
 =\frac{R!}{\prod_i x_i!}
       \prod_i\left(\frac{\operatorname{mult}_i}{D_{\mu_i}}
                         T_b^{\mu_i}(p)\right).
\]
Using \eqref{s:hybrid-local-ratio}, the remaining binomial and
position factors, including \eqref{s:hybrid-global-marker-multiplicity},
are bounded by
\[
 \binom nR(e/a)^R
       \frac{R!}{\prod_i x_i!}\prod_i\binom b{x_i}^{-1}
 \le \binom nR R!(e/a)^R(e/b)^R
 \le e^{2R}.
\]
Thus the normalized trace endpoint obeys
\begin{equation}\label{rec:hybrid-marked-trace}
 \begin{split}
 \log\left(\frac{D_\lambda}{L_\lambda}
           \Tr_{\mathcal H}(B_\beta^{2p}A_\alpha^{2p})\right)
 \le{}&\sum_{\rm children}\log T
       +G-G_H-G_V+\sum_iK_i\\
     &+O(R+(a+b)D_0^2\log^2(en)).
 \end{split}
\end{equation}

\paragraph{Charging the child diagrams.}
We next compare the child diagram budgets with the chain entropies.
Keep the maximizing shapes $\rho_i$ chosen in
\eqref{s:hybrid-local-ratio}. Call a horizontal chain $(i,h)$ good
if $l_{ih}\ge b^{1-\alpha_0}$, where
$\alpha_0=v_*/(4c_*)$, and define the number of good extensions by
\[
 e_{\rm good}=\sum_{\substack{i,h\\ l_{ih}\ge b^{1-\alpha_0}}}e_{ih}.
\]
This counts horizontal extension boxes only, so
$0\le e_{\rm good}\le R$.

The $l^0_{ih}$ unmarked boxes in this chain each see a full row or
column length in $\mu_i$ at least $l_{ih}\ge l^0_{ih}$.
Their contribution to $\mathcal D_{\log_2 b}(\mu_i)$ is at most
$c_{\log_2 b}\sum_hl^0_{ih}\log(b/l^0_{ih})$. By
\eqref{s:hybrid-local-chain-entropy} and
$M_i\log(b/M_i)\le x_i$, this is
$c_{\log_2 b}G_i+O(x_i+D_0^2\log(en))$.
Charge the remaining $x_i$ boxes by the cap
$v_{\log_2 b}\log b$. For a good extension box the actual charge
is at most
$c_{\log_2 b}\alpha_0\log b\le(v_*/2)\log b$,
since $c_{\log_2 b}\le2c_*$. The cap is at least $v_*\log b$,
so each such box saves at least $(v_*/2)\log b$.
In columns choose any compatible unmarked shape and make the same
estimate, discarding its extension savings.

Set
$\bar c=\max(c_{\lfloor d/2\rfloor},c_{\lceil d/2\rceil})$ and
$\bar v=\max(v_{\lfloor d/2\rfloor},v_{\lceil d/2\rceil})$.
The marker counts sum to $R$ on each axis and
$\log a+\log b=\log n$. Therefore
\begin{equation}\label{rec:hybrid-child-diagram}
 \sum_{\rm children}\mathcal D_{d_i}(\mu_i)
 \le\bar c(G_H+G_V)+\bar vR\log n
   -(v_*/2)e_{\rm good}\log m
   +O(R+(a+b)D_0^2\log(en)).
\end{equation}

It remains to bound the $K_i$ costs of the extensions.
We may restrict to profile pairs with nonzero angle, since the other
pairs contribute zero to \eqref{rec:hybrid-two-strip}.
We first show that the total length $U_{\rm bad}$ of bad horizontal
chains is at most $2k$, for all sufficiently large $d$.
If $a_1=0$, then $\lambda_1<n^{1-\gamma_d}\le n/2$ and hence
$k\ge n/2$; the total of all horizontal chain lengths is at most
$n\le2k$.
If $a_1>0$, the nonzero space $E_\alpha S\mathcal H$ contains the
global compact type $(\eta,\zeta)$ inside the tensor product of
the prescribed row types. Its highest weight has first even
coordinate $\eta_1=\lambda_1$. Weights add under the diagonal color
action, and in local type $\eta^{(i)}$ every first weight coordinate is
at most $\eta^{(i)}_1$. Thus
\[
 \lambda_1=\eta_1\le\sum_i\eta^{(i)}_1
                  \le\sum_i\mu_{i,1},\qquad
 \sum_i(b-\mu_{i,1})\le k.
\]
The second inequality uses $\eta^{(i)}\subseteq\rho_i\subseteq\mu_i$.
All chains other than a row's first chain lie below that first row.
If the first chain is itself bad, its length
$\mu_{i,1}<b^{1-\alpha_0}\le b/2$ is at most the same row's
minority $b-\mu_{i,1}$. Summing these two contributions gives
$U_{\rm bad}\le2k$.

For any collection of chains with total length $U$ and total
extension count $z>0$, Vandermonde's identity and the binomial bound
give
\[
 \sum_h\log\binom{l_h}{e_h}
 \le\log\binom Uz\le z\log(eU/z).
\]
For bad chains $U\le2k$ and $z\le R$. Using
$\log x\le\log(d+1)+x/(d+1)$ yields the explicit bound
\[
 K_{\rm bad}\le
 R\log(d+1)+\frac{2ek}{d+1}
 =O(R\log(d+1)+k/d).
\]
Zero extension counts contribute zero.
For good chains $U\le n$ and $z=e_{\rm good}$.
If $e_{\rm good}\ge k n^{-\epsilon}$, the dense hypothesis
$k\ge n^{1-\epsilon}$ gives
\[
 K_{\rm good}\le
 e_{\rm good}\log(en/e_{\rm good})
 \le e_{\rm good}(1+2\epsilon\log n)
 \le (v_*/2)e_{\rm good}\log m
\]
for sufficiently large $d$. The final inequality follows from
$2\epsilon<v_*/4$ and
$\log m\ge(\log n-\log2)/2$.
If $0<e_{\rm good}<k n^{-\epsilon}$, monotonicity of
$z\log(en/z)$ on $[0,n]$ instead gives
$K_{\rm good}=O(k n^{-\epsilon}\log(en))$.
Consequently the $K_i$ terms and the negative term in
\eqref{rec:hybrid-child-diagram} together cost at most
\begin{equation}\label{s:hybrid-extension-cost}
 O\!\left(R\log(d+1)+k/d+k n^{-\epsilon}\log(en)\right).
\end{equation}
All thresholds here depend only on the displayed fixed parameters,
not on the eventual exponent $p$.

The color angle supplies the remaining saving:
\begin{equation}\label{rec:hybrid-color-angle}
 \log w\le-\tfrac12(G-G_H-G_V)
                    +O(R+(a+b)D_0^2\log n),\qquad w\le1.
\end{equation}
If $M=0$, there are no unmarked symbols, $G=G_H=G_V=0$, and
\eqref{rec:hybrid-color-angle} follows from $w\le1$.  Assume $M>0$
for the density and deformation argument that follows.
We first pass from the full local color projections to weight
projections. Let $Q_H,Q_V$ be the products of the specified local
color projections on the rows and columns. Write the additional
permutation-type projections as $E^{\rm perm},F^{\rm perm}$, so
$E_\alpha=Q_HE^{\rm perm}$ and $F_\beta=F^{\rm perm}Q_V$.
They commute within their respective lines and with $S$. Hence
\[
 F_\beta E_\alpha S
   =F^{\rm perm}(Q_V S Q_H)E^{\rm perm},
 \qquad w\le\|Q_V S Q_H\|.
\]
All three color projections preserve marker positions, so this norm
is the maximum of the norms on fixed marker-position fibers.

On one such fiber, let $Q_\tau$ be a local color-type projection,
and let $d_\tau$ be its compact carrier dimension. In every
multiplicity copy choose the highest-weight line and let
$P_\tau^{\rm hw}$ project onto their sum. Normalized Haar averaging
on $K=U(a_1)\times U(b_1)$ gives, with $\rho(g)$ denoting the
local tensor color action,
\[
 Q_\tau=d_\tau\int_K
       \rho(g)P_\tau^{\rm hw}\rho(g)^*\,dg.
\]
The highest-weight lines lie in the torus weight space whose symbol
counts are the parts of $\tau$. If $Z_\tau$ projects onto that
whole weight space, then $P_\tau^{\rm hw}\preceq Z_\tau$. Tensor
these positive inequalities over the rows and over the columns and
apply Lemma~\ref{ten:positive-mixtures}, with middle operator $S$.
It follows that
\[
 \|Q_V S Q_H\|
 \le\left(\prod_{\rm lines}d_\tau\right)^{1/2}
      \sup_{\mathbf g,\mathbf h}
       \|Z_V(\mathbf h)S Z_H(\mathbf g)\|.
\]
Here the rotations are independent on each line. The logarithm of
the prefactor is $O((a+b)D_0^2\log n)$ by the Weyl dimension
formula. This argument applies with either parity dimension zero
by omitting that factor of $K$; it uses no commutation between row
and column projections.

Fix the line rotations in this last supremum. If a row has
$M_i>0$ unmarked sites,
let $\sigma_i$ be the parity-block diagonal density whose entries
are its required symbol counts divided by $M_i$, in the chosen
rotated parity basis. If $M_i=0$,
choose any parity-preserving density of trace one instead; that row
has no tensor factors in the fiber.  Define column densities
$\tau_j$ in the same way, using an arbitrary parity-block trace-one
density when its local unmarked count is zero.
Thus every line density has trace one, including empty lines.
On the row weight space $Z_H(\mathbf g)$ the tensor multiplier $\Sigma$ of the
$\sigma_i^{1/2}$ acts by $e^{-G_H/2}$, and the analogous column
multiplier $\mathcal T$ acts by $e^{-G_V/2}$. In particular,
\[
 Z_V(\mathbf h)S Z_H(\mathbf g)
 =e^{(G_H+G_V)/2}
   Z_V(\mathbf h)\mathcal T S\Sigma Z_H(\mathbf g).
\]
Thus the rotated weight-space overlap is at most
$e^{(G_H+G_V)/2}\|\mathcal T S\Sigma\|$.
For the global compact projection use the deformed formula
with $\bar\tau=b^{-1}\sum_j\tau_j$,
$\bar\sigma=a^{-1}\sum_i\sigma_i$, and a real scalar regularizer
$\varepsilon_{\rm reg}>0$.  For an invertible parity-preserving block
matrix $A$, the polynomial tensor action
$\rho(A)=A^{\otimes M}$ acts on the whole color tensor fiber, while
$\pi(A)$ acts on the selected carrier $\mathcal U_{\eta,\zeta}$.
The restriction of $S$ to this fiber is the full compact
$\pi$-isotypic projection. Put
$A_0=(\bar\tau+\varepsilon_{\rm reg}I)^{-1/2}$.
With $dU$ denoting normalized Haar measure
on $K$, apply Lemma~\ref{s:compact-coefficient-extraction} with carrier
matrix $\pi(A_0)^{-1}$, and multiply its tensor action on the left by
$\rho(A_0)$. The formula is
\[
 \begin{gathered}
 S=(\dim\pi)\int_K
 \Tr(\pi(U)^*\pi(A_0)^{-1})\,\rho(A_0U)\,dU.
 \end{gathered}
\]
To bound the trace coefficient, put
\[
 H=\bar\tau+\varepsilon_{\rm reg}I,\qquad
 t=\operatorname{Tr}H=1+\varepsilon_{\rm reg}(a_1+b_1),\qquad
 D=H/t.
\]
Since $M>0$, the total color dimension $a_1+b_1$ is positive,
and $D$ is a trace-one parity-preserving density. Apply
\eqref{eq:signed-carrier-density-upper} with parity dimensions
$a_1,b_1$ and partitions $\eta,\zeta$. Their total degree is $M$
and their entropy is $G$, so $\|\pi(D)\|\le e^{-G}$.
Polynomial homogeneity and positivity give
\[
 \begin{aligned}
 \pi(A_0)^{-1}&=\pi(H^{1/2})=t^{M/2}\pi(D)^{1/2},\\
 \|\pi(A_0)^{-1}\|
       &\le(1+\varepsilon_{\rm reg}(a_1+b_1))^{M/2}e^{-G/2}.
 \end{aligned}
\]
Here the square-root identity follows by diagonalizing $D$ within
its parity blocks. An absent block contributes the trivial carrier
factor and degree zero. Since $\pi(U)$ is unitary, the absolute
trace coefficient is at most $\dim\pi$ times this norm, in addition
to the $\dim\pi$ outside the integral.
For the site factor $\tau_j^{1/2}A_0U\sigma_i^{1/2}$, the squared
Hilbert--Schmidt norms have full-grid mean
\[
 \Tr(A_0\bar\tau A_0\,U\bar\sigma U^*)\le1,
\]
because $A_0\bar\tau A_0\le I$ and $\Tr\bar\sigma=1$.
Thus on the $M$ unmarked sites the product of the site norms is at most
$(n/M)^{M/2}\le e^{R/2}$.  Let $\varepsilon_{\rm reg}$ decrease to zero.
Including the carrier dimensions proves
\eqref{rec:hybrid-color-angle}.

\begin{proof}[Completion of Theorem~\ref{rec:hybrid-main}]
The number $J$ of profile pairs in this application satisfies
\[
 \log J=O\!\left(n^{3/4}+(a+b)D_0\log(en)\right).
\]
Indeed the permutation partitions contribute
$O(a\sqrt b+b\sqrt a)=O(n^{3/4})$ to this logarithm. Every local
compact type is specified by at most $a_1+b_1$ nonnegative parts,
each at most $n$, giving the second term over all lines. Marker
counts are fixed by the local color degrees and introduce no
additional profile choices.

Set $\Delta G=G-G_H-G_V$. Combining
\eqref{rec:hybrid-marked-trace},
\eqref{rec:hybrid-child-diagram},
\eqref{s:hybrid-extension-cost}, and
\eqref{rec:hybrid-color-angle} in
\eqref{rec:hybrid-two-strip}, then subtracting
\eqref{rec:hybrid-arm-weight}, gives
\[
 \begin{split}
 \log T_n^\lambda(p)-\mathcal D_d(\lambda)
 \le{}&(1-\bar c)\Delta G-(p_0-1)(\Delta G)_+\\
 &-(c_d-\bar c)G-(v_d-\bar v)R\log n\\
 &+O_{p_0}\!\left(R\log(d+1)+k/d\right).
 \end{split}
\]
To obtain this error, use $k\ge n^{1-\epsilon}$,
$D_0=O(n^{0.01})$, and $a+b=O(\sqrt n)$. The terms
$(a+b)D_0^2\log^2(en)$, $\log J$, and
$kn^{-\epsilon}\log(en)$ are all $o(k/d)$.
The angle is bounded both by its exponential estimate and by one;
therefore its main term is $-(p_0-1)(\Delta G)_+$, with the stated
$O_{p_0}$ error. No factor in this error depends on the final
choice of $p$.

The first line of the display is nonpositive. If $\Delta G\le0$,
use $1-\bar c>0$; if $\Delta G>0$, use
$1-\bar c-(p_0-1)<0$.
Let $Q=G+R\log n$. By \eqref{rec:hybrid-arm-dimension},
$Q\ge k/2$, and
\[
 R\log(d+1)+k/d
 \le C Q\,\frac{\log(d+1)}d.
\]
The balanced parameter increments satisfy
$c_d-\bar c=v_d-\bar v\ge c d^{-1/3}$.
Their negative contribution is at most $-cd^{-1/3}Q$ and absorbs
the error for sufficiently large $d$. This proves the dense
diagram induction.

For its comparison with dimension, put
$Q_\gamma=G+\gamma_dR\log n$. The fixed parameters give
$\gamma_d\ge10/10001>0$, and $\gamma_d<1$, so
$Q_\gamma\ge cQ\ge ck$.
Moreover $R=o(Q_\gamma)$ uniformly, because
$Q_\gamma\ge\gamma_dR\log n$, and
$D_0^2\log^2(en)=o(Q_\gamma)$ in the dense range. Thus
\eqref{rec:hybrid-arm-weight} and
\eqref{rec:hybrid-arm-dimension} give
\[
 \mathcal D_d(\lambda)=c_dQ_\gamma+o(Q_\gamma),
 \qquad
 \log D_\lambda\ge Q_\gamma-o(Q_\gamma).
\]
The fixed parameter bounds, in particular $c_d<0.011$, prove
$\mathcal D_d(\lambda)\le\frac14\log D_\lambda$ and
$\log D_\lambda\ge ck$ for all sufficiently large $d$.  This
comparison used only the stated height bound and the diagram estimates,
so it also applies to zero blocks in the theorem's domain.
The other levels were covered by the level induction. All thresholds
are chosen before the finite-base exponent $p$.

Choose the finite-base exponent using \eqref{rec:hybrid-explicit-base}.

In the dense range $\|B_n\|^{2p}\le D_\lambda^{-3/4}$ and
$T_n^\lambda(p)\le D_\lambda^{1/4}$.  Thus
$T_n^\lambda(4p)\le D_\lambda^{-2}$.
In the other range $T_n^\lambda(p)\le n^{-k}$, whence
$\|B_n\|^{2p}\le D_\lambda^{-1}n^{-k}$ and
$T_n^\lambda(4p)\le D_\lambda^{-3}n^{-4k}$.
Consequently exactly $4p$ sweeps suffice:
Lemma~\ref{lem:schatten}, with $r=s=4p$, gives
\begin{equation}\label{rec:hybrid-exact-completion}
 \|B_n^{4p}|_{\mathbf1^\perp}\|_{\HS,\mathrm{reg}}^2
 \le\sum_{\lambda\ne(n)}T_n^\lambda(4p)=o(1).
\end{equation}
The dense sum tends to zero because $k\ge n^{1-\epsilon}$ and
dimensions dominate the partition count; the remaining sum is
bounded by $\sum_{k\ge1}e^{O(\sqrt k)}n^{-4k}$.
Thus $4pd$ physical shuffles suffice for all sufficiently large $d$.
\end{proof}
